
\section{Additional Model Training Details}

\subsection{Text encoders}
As described in~\cref{sec:t2v_text_encoder} we use text embeddings from three text encoders.
We provide details on the text encoders and how they are used for visual-text generation.

\noindent \textbf{Long-prompt MetaCLIP training.}
We train our own CLIP-style~\citep{xu2023demystifying} model that can process long text prompts up to $256$ text tokens.
We utilized synthetic image captions generated by an image captioning model~\citep{llama3} to finetune the MetaCLIP~\citep{xu2023demystifying} text encoder.
We expanded the position embedding in the text encoder to $256$ tokens, allowing it to handle longer input sequences.
However, we found that finetuning all parameters in the text encoder led to rapid overfitting and suboptimal text encoding performance.
We hypothesize that this is due to the uniform text style in the synthetic long image captions.
To address this issue, we froze both the image and text encoders and finetuned the position embedding and all biases in the text encoder using a mix of long synthetic, short synthetic, and human-annotated captions.

\par \noindent \textbf{Visual text generation.}
We implemented several key modifications to the input and prompt rewriting, model architecture, and data processing steps to enable visual text generation.
Specifically, our approach consists of three main components.
First, we assume that text enclosed within quotation marks (\textbf{`` ''}) in the input text prompt needs to be generated as visual-text.
We also employ prompt rewriting (see~\cref{sec:t2v_prompt_rewrite}) to automatically identify such text at inference time.
Second, we use the character-level ByT5 encoder for encoding the text within the quotation marks.
Finally, we ensure that our \pretraining data has a good mix of visual text that covers 10-50\% of the image area by leveraging OCR detection models.





\subsection{Model scaling and training efficiency}
\label{app:t2v_impl_scaling_additional}

For memory capacity constrained models like \OursVideo with a context length of 73K, sub-optimal performance may be achieved through memory-saving techniques like activation checkpointing (AC).
AC reduces peak memory demand by trading off FLOPs and memory, which can be sometimes be more effective than fully-optimized parallelisms due to physical training system constraints, \eg, the FLOPs/sec and HBM capacity of each GPU, and the intra- and inter- node GPU-GPU interconnect bandwidths.

\noindent \textbf{Overlapping communication and computation.}

While the parallelism techniques mentioned in \cref{sec:t2v_impl_scaling}—which aim to partition training FLOP and memory demands across GPUs at the cost of added communication—are successful at enabling the 0-to-1 ability to train such large sequence transformer models, their direct implementation and composition come with overheads and inefficiencies with respect to memory and communication. As a result, under such overheads optimal \textit{realized} performance for a model which is memory capacity constrained, as-is \OursVideo with a context length of 73K, may actually be achieved through the use of other memory-saving techniques in addition to, or even in place of, the above parallelisms.

Given that AC can never provide strong scaling due to strictly increasing the amount of work needed to compute a training step, we focused our efforts on improving the scaling characteristics of model parallelism techniques. Specifically, we aimed to have a final model parallelism implementation which achieved as close to strong scaling as possible for both: 1) activation memory size (to reduce the usage of AC), and 2) forward/backward step time.

As the four existing model parallel techniques discussed in~\cref{sec:t2v_impl_scaling} (TP, SP, CP, FSDP) all already achieve strong theoretical FLOP scaling, we began by: 1) determining the scaling characteristics of the activation sizes of our model under these techniques as well as 2) building an analytical framework to model the change and interdependence of the compute and communication times of our model execution. Using both, we: 1) identified all occurrences of duplicated activations, and 2) identified which inter-GPU communications are required as well as exposed.

From this we designed and implemented a model parallel solution for the \OursVideo backbone based upon the principles of TP, SP, and CP, which: 1) achieves strong activation memory scaling, and 2) minimizes exposed communication time for our training cluster. This enabled us to achieve close-to-strong scaling even with model parallelism widths spanning multiple nodes. We achieved such performance through a custom implementation defining the execution of both the forward and backward paths of the \OursVideo backbone block. Forgoing the use and ease of chaining smaller autograd operators together allowed us to precisely partition, control, and overlap compute and communication execution. This also enabled us to transparently apply techniques such as selective recomputation without any performance loss.

Our implementation is fully written in PyTorch and compiles into CUDAGraphs supporting the execution of variable sized inputs at the start of every training or inference job, dynamically based on the specified model and system configurations.

\noindent \textbf{Sharding plan generation and selection.}
We expanded the performance modeling tools developed above to further estimate the memory utilization and latency of end-to-end model execution (\eg, backbone blocks, text encoders, \taeShort). We then utilized this to generate multiple sharding and parallelism plans, which exhibit similar theoretical and estimated latencies, and are deemed valid as they are estimated to fit within the available GPU High-Bandwidth Memory (HBM).

This development: 1) enabled us to generate sharding plans in which different components and stages of the end-to-end model execution can be sharded by different strategies; and 2) allowed us to empirically identify training parallelism configurations which have an approximately neutral batch-size to step time scaling relationship.

The ability to accomplish the latter was important for the successful training of \OursVideo due to the relationship of model parallel size to the global batch size of its corresponding training step. Specifically, both the number of GPUs over which parameters are disjointly sharded (TP) and the number of GPUs over which the input sequence of a single sample are disjointly sharded (SP, CP), proportionally reduce the effective global batch size—and impact the wall latency of—the corresponding training step. Identifying groups of sharding plans which have neutral batch-size to step time scaling allows us to effectively scale the number of optimizer steps, while holding the number of GPUs and total training data size constant. Although the training data throughput and end-to-end efficiency of such groups of sharding plans are similar, the number of optimizer steps taken to process varying amounts of training data across various stages of training can significantly impact the final model's quality.

The final sharding plan used during the the most expensive stage of \OursVideo training, processing 768px video inputs with a per-sample token sequence length of 73K, was the following:
\begin{itemize}
  \item \textbf{Text encoders.} Due to their relatively small size and weights being frozen, ByT5 and Long-Prompt MetaCLIP were replicated on all GPUs. UL2 however has significantly more parameters and memory overhead, yet it is still relatively small in terms of end-to-end execution latency, and was sharded FSDP-only across the DP-group of each TP-rank.
  \item \textbf{\taeShort.} Although containing a relatively small number of frozen parameters, the size of intermediate activations of the \taeShort can become prohibitively expensive as the input size grows. Furthermore, unlike the text encoders, the latency of the \taeShort is non-trivial with respect to the end-to-end step time, and unlike the backbone, it is non-trivial to efficiently partition and model parallelize the \taeShort's execution. These limitations resulted in us performing a data pre-processing step where the latents for high resolution video inputs were pre-computed and cached prior to their ingestion in the \OursVideo backbone training pipeline.
  \item \textbf{\OursVideo} \textbf{backbone.} While a transformer block at its core, the \OursVideo backbone has additional learned components, such as factorized positional embeddings and per-context embedders, each with not only their own memory and compute requirements but also their own input and output activation connections. This results in the backbone parameter sharding and input and output activation and gradient flow changing as the model moves through its different stages. The final backbone contained interconnected sections sharded: FSDP-only (\eg, patchifier), FSDP+TP (\eg, context embedders), FSDP+TP+SP (\eg, cross-attention), and FSDP+TP+SP+CP (\eg, self-attention).
\end{itemize}



